% Book 3, Year 1, chapter 11 -- induction, recursion and termination.
% Provenance (never printed): sources/book3/BRIEFS.md, chapter 11.
\chapter{Induction, Recursion and Termination}\label{ch:b1:induction-and-termination}

A function reverses a list; applied twice, it should give the list back. A program can check this on three thousand
lists in a fraction of a second, and it does below, but lists have no largest length: a check, however long, says
nothing about the lists it did not try. This chapter gives the two tools that cover every case. \omterm{def:b1:induction-and-termination:structural}{Structural induction}
proves a property of every list, tree or formula by following the rules that build them. \omterm{def:b1:induction-and-termination:well-founded}{Well-founded orders} prove
that a loop or a recursive function stops, by attaching to each state a quantity that cannot decrease forever.

\begin{recall}
To prove a property $P(n)$ for every natural number $n$ by induction, prove $P(0)$, and prove $P(n + 1)$ from $P(n)$
(or from all of $P(0), \dots, P(n)$: strong induction). An order relation is reflexive, antisymmetric and transitive.
A recursive function calls itself on smaller arguments and stops on base cases.
\end{recall}

\section{Inductively defined sets}\label{sec:b1:induction-and-termination:sets}

\begin{definition}[Inductively defined set]\label{def:b1:induction-and-termination:inductive}
A set is \emph{inductively defined}\index{inductively defined set} by rules when it is the smallest set closed under
them: \emph{base rules} put given elements in the set, and \emph{construction rules} build an element from elements
already in it. Every element is obtained by finitely many applications of the rules.
\end{definition}

The natural numbers are defined by the base rule ``0'' and the rule ``if $n$, then $n + 1$''; the lists of the
elements of a set $A$ by ``the empty list'' and ``if $x \in A$ and $\ell$ is a list, then $x$ followed by $\ell$'';
the \omterm{def:b1:propositional-logic:formula}{propositional formulas} of \cref{ch:b1:propositional-logic} by their variables and \omterm{def:b1:propositional-logic:formula}{connectives}.

\begin{definition}[Rooted tree, root, child, leaf, internal node, subtree]\label{def:b1:induction-and-termination:rooted-tree}
A \emph{rooted tree}\index{rooted tree} is a node, its \emph{root}\index{root (of a tree)}, holding a key and a finite
sequence of rooted trees, the subtrees of the root; the roots of these subtrees are the \emph{children}\index{child (of a
node)} of the root, which is their \emph{parent}. A node without children is a \emph{leaf}\index{leaf (of a tree)};
the others are \emph{internal nodes}\index{internal node}. The \emph{subtrees}\index{subtree} of a tree are the tree
and the subtrees of its subtrees.
\end{definition}

\begin{definition}[Binary tree, full binary tree]\label{def:b1:induction-and-termination:binary-tree}
The set of \emph{binary trees}\index{binary tree} is defined inductively: the \emph{empty tree} is a binary tree, and if
$\ell$ and $r$ are binary trees and $k$ a key, the node $(\ell, k, r)$, with left \omterm{def:b1:induction-and-termination:rooted-tree}{subtree} $\ell$ and right \omterm{def:b1:induction-and-termination:rooted-tree}{subtree} $r$,
is a binary tree. A binary tree is \emph{full}\index{full binary tree} if every node has either two non-empty \omterm{def:b1:induction-and-termination:rooted-tree}{subtrees}
or two empty ones.
\end{definition}

\begin{definition}[Depth, size, height]\label{def:b1:induction-and-termination:height}
The \emph{depth}\index{depth (of a node)} of a node is the number of links from the \omterm{def:b1:induction-and-termination:rooted-tree}{root} to it. The
\emph{size}\index{size (of a tree)} of a \omterm{def:b1:induction-and-termination:binary-tree}{binary tree} is its number of nodes, 0 for the empty tree. Its
\emph{height}\index{height (of a tree)} is the largest depth of its nodes, and $-1$ for the empty tree: by induction,
$h(\text{empty}) = -1$ and $h((\ell, k, r)) = 1 + \max(h(\ell), h(r))$.
\end{definition}

\omcode{code/bachelor-1/11-induction-and-termination/listings/inductive.py}{6}{35}{\omterm{def:b1:induction-and-termination:binary-tree}{Binary trees}, and every tree of
\omterm{def:b1:induction-and-termination:height}{height} at most $h$}\label{lst:b1:induction-and-termination:inductive}

\omoutput{code/bachelor-1/11-induction-and-termination/outputs/inductive.txt}{A tree of four nodes, the empty tree,
and the trees produced by the rules up to \omterm{def:b1:induction-and-termination:height}{height} 3.}\label{out:b1:induction-and-termination:inductive}

\code{shapes} applies the rules literally: the trees of \omterm{def:b1:induction-and-termination:height}{height} at most $h$ are the empty tree and the nodes built from
two trees of \omterm{def:b1:induction-and-termination:height}{height} at most $h - 1$, so their number $a_h$ satisfies $a_{-1} = 1$ and $a_h = 1 + a_{h-1}^2$. The
convention $-1$ for the empty tree makes this rule, and the bounds below, hold without exceptions.

\begin{omfigure}
\begin{tikzpicture}[n/.style={circle, draw, fill=omCode!10, minimum size=4mm, inner sep=0pt},
                    e/.style={rectangle, draw, dashed, minimum size=2.5mm, inner sep=0pt}, level distance=7mm,
                    sibling distance=7mm]
\node[draw, rounded corners=2pt, align=left, font=\scriptsize] at (-1.2,0.2) {rule 1: the empty tree
\tikz\node[e]{}; \\ rule 2: from trees $\ell$, $r$, the node\\[1mm] \hspace*{8mm}\tikz[level distance=6mm, sibling
distance=7mm]\node[n]{} child {node[font=\scriptsize] {$\ell$}} child {node[font=\scriptsize] {$r$}};};
\node[font=\scriptsize\bfseries] at (6.2,1.0) {the 5 binary trees of height at most 1};
\node[e] at (3.6,0) {};
\node[n] at (4.8,0) {} child {node[e] {}} child {node[e] {}};
\node[n] at (6.2,0) {} child {node[n] {} child {node[e] {}} child {node[e] {}}} child {node[e] {}};
\node[n] at (7.6,0) {} child {node[e] {}} child {node[n] {} child {node[e] {}} child {node[e] {}}};
\node[n] at (9.0,0) {} child {node[n] {} child {node[e] {}} child {node[e] {}}} child {node[n] {} child {node[e] {}}
  child {node[e] {}}};
\end{tikzpicture}
\par\smallskip{\small The two rules that generate \omterm{def:b1:induction-and-termination:binary-tree}{binary trees}, and what they produce from trees of \omterm{def:b1:induction-and-termination:height}{height} at most 0
(the empty tree, dashed, and the single node). \omterm{def:b1:induction-and-termination:height}{Height} $-1$, 0, 1, 1 and 1; the last tree is \omterm{def:b1:induction-and-termination:binary-tree}{full}.}
\end{omfigure}

\section{Structural induction}\label{sec:b1:induction-and-termination:structural}

\begin{definition}[Structural induction]\label{def:b1:induction-and-termination:structural}
To prove a property $P$ of every element of an \omterm{def:b1:induction-and-termination:inductive}{inductively defined set} by \emph{structural
induction}\index{structural induction}, prove that $P$ holds for the elements given by the base rules, and that each
construction rule, applied to elements satisfying $P$, builds an element satisfying $P$.
\end{definition}

\begin{proposition}[Structural induction is sound]\label{prop:b1:induction-and-termination:structural}
If the two parts of a \omterm{def:b1:induction-and-termination:structural}{structural induction} are proved, then $P$ holds for every element of the set.
\end{proposition}

\begin{proof}
Every element is built by finitely many applications of the rules; call its \emph{rank} the least number of rule
applications in such a construction. We prove ``$P$ holds for every element of rank at most $n$'' by induction on
$n \geq 1$. An element of rank 1 comes from a base rule: $P$ holds. An element of rank $n + 1$ comes from a
construction rule applied to elements of smaller rank, which satisfy $P$ by the induction hypothesis, so it satisfies
$P$. (For \omterm{def:b1:induction-and-termination:binary-tree}{binary trees}, an induction on the \omterm{def:b1:induction-and-termination:height}{height} does the same work.)
\end{proof}

\begin{proposition}[Size, height and leaves of a binary tree]\label{prop:b1:induction-and-termination:bounds}
A \omterm{def:b1:induction-and-termination:binary-tree}{binary tree} of \omterm{def:b1:induction-and-termination:height}{height} $h$ has a \omterm{def:b1:induction-and-termination:height}{size} $n$ with $h + 1 \leq n \leq 2^{h + 1} - 1$. A \omterm{def:b1:induction-and-termination:binary-tree}{full} non-empty \omterm{def:b1:induction-and-termination:binary-tree}{binary tree} has one
more \omterm{def:b1:induction-and-termination:rooted-tree}{leaf} than \omterm{def:b1:induction-and-termination:rooted-tree}{internal nodes}.
\end{proposition}

\begin{proof}
By \omterm{def:b1:induction-and-termination:structural}{structural induction}. The empty tree has $h = -1$ and $n = 0 = h + 1 = 2^0 - 1$. For a node with \omterm{def:b1:induction-and-termination:rooted-tree}{subtrees} of
\omterm{def:b1:induction-and-termination:height}{heights} $h_\ell$, $h_r$ and \omterm{def:b1:induction-and-termination:height}{sizes} $n_\ell$, $n_r$ satisfying the bounds, $h = 1 + \max(h_\ell, h_r)$ and $n = 1 + n_\ell +
n_r$. The higher \omterm{def:b1:induction-and-termination:rooted-tree}{subtree} alone has at least $\max(h_\ell, h_r) + 1 = h$ nodes, so $n \geq h + 1$; and $n \leq 1 +
2(2^{h} - 1) = 2^{h + 1} - 1$, since both \omterm{def:b1:induction-and-termination:height}{heights} are at most $h - 1$. For \omterm{def:b1:induction-and-termination:binary-tree}{full} trees: a single node has one \omterm{def:b1:induction-and-termination:rooted-tree}{leaf} and no
\omterm{def:b1:induction-and-termination:rooted-tree}{internal node}. A node with two non-empty \omterm{def:b1:induction-and-termination:binary-tree}{full} \omterm{def:b1:induction-and-termination:rooted-tree}{subtrees} having $f_\ell = i_\ell + 1$ and $f_r = i_r + 1$ leaves has
$f_\ell + f_r = i_\ell + i_r + 2$ leaves and $i_\ell + i_r + 1$ \omterm{def:b1:induction-and-termination:rooted-tree}{internal nodes}, itself included.
\end{proof}

\begin{omfigure}
\begin{tikzpicture}[n/.style={circle, draw, fill=omDef!15, minimum size=4.5mm, inner sep=0pt},
                    l/.style={circle, draw, fill=omCode!10, minimum size=4.5mm, inner sep=0pt}, level distance=8mm,
                    level 1/.style={sibling distance=26mm}, level 2/.style={sibling distance=13mm},
                    level 3/.style={sibling distance=7mm}]
\node[n] {} child {node[n] {} child {node[l] {}} child {node[n] {} child {node[l] {}} child {node[l] {}}}}
            child {node[n] {} child {node[l] {}} child {node[l] {}}};
\draw[dashed, omThm, rounded corners] (-2.4,-0.5) rectangle (0.05,-2.85);
\draw[dashed, omThm, rounded corners] (0.25,-0.5) rectangle (2.4,-2.0);
\node[font=\scriptsize, align=left, anchor=west] at (2.8,-0.4) {left subtree: 3 leaves, 2 internal\\right subtree:
2 leaves, 1 internal\\the whole: $3 + 2 = 5$ leaves,\\$2 + 1 + 1 = 4$ internal (the root added)};
\end{tikzpicture}
\par\smallskip{\small The induction step for \omterm{def:b1:induction-and-termination:binary-tree}{full} trees: each \omterm{def:b1:induction-and-termination:rooted-tree}{subtree} has one more \omterm{def:b1:induction-and-termination:rooted-tree}{leaf} than \omterm{def:b1:induction-and-termination:rooted-tree}{internal nodes}, and the
\omterm{def:b1:induction-and-termination:rooted-tree}{root} adds one \omterm{def:b1:induction-and-termination:rooted-tree}{internal node} and no \omterm{def:b1:induction-and-termination:rooted-tree}{leaf}, so the whole tree keeps the difference of one. Leaves are pale, \omterm{def:b1:induction-and-termination:rooted-tree}{internal nodes}
shaded.}
\end{omfigure}

Lists are proved about in the same way. With the recursive definition $\mathit{rev}([\,]) = [\,]$ and
$\mathit{rev}([x] + \ell) = \mathit{rev}(\ell) + [x]$, first prove $\mathit{rev}(\ell + m) = \mathit{rev}(m) +
\mathit{rev}(\ell)$ by induction on $\ell$: for $\ell = [\,]$ both sides are $\mathit{rev}(m)$; and
$\mathit{rev}([x] + \ell + m) = \mathit{rev}(\ell + m) + [x] = \mathit{rev}(m) + \mathit{rev}(\ell) + [x] =
\mathit{rev}(m) + \mathit{rev}([x] + \ell)$. Then $\mathit{rev}(\mathit{rev}(\ell)) = \ell$ by induction on $\ell$:
$\mathit{rev}(\mathit{rev}([x] + \ell)) = \mathit{rev}(\mathit{rev}(\ell) + [x]) = \mathit{rev}([x]) +
\mathit{rev}(\mathit{rev}(\ell)) = [x] + \ell$.

\omcode{code/bachelor-1/11-induction-and-termination/listings/reverse_twice.py}{6}{21}{Reverse, and a check on 3280
lists}\label{lst:b1:induction-and-termination:reverse}

\omoutput{code/bachelor-1/11-induction-and-termination/outputs/reverse_twice.txt}{The check passes; the proof above
covers the lists it did not try.}\label{out:b1:induction-and-termination:reverse}

The first lemma was needed because the property of $\mathit{rev}(\mathit{rev}(\ell))$ alone does not carry through the
induction step: proving a stronger or more general statement is often what makes an induction go through.

\section{Well-founded orders}\label{sec:b1:induction-and-termination:orders}

\begin{definition}[Strict order, minimal element, immediate predecessor]\label{def:b1:induction-and-termination:strict-order}
A \emph{strict order}\index{strict order} $<$ on a set is irreflexive ($x < x$ never holds) and transitive; from an
order $\leq$ one gets the strict order $x < y$ when $x \leq y$ and $x \neq y$. An element $m$ of a subset $S$ is
\emph{minimal}\index{minimal element} in $S$ if no element of $S$ is smaller than $m$. An \emph{immediate
predecessor}\index{immediate predecessor} of $y$ is an $x < y$ with no $z$ such that $x < z < y$.
\end{definition}

\begin{definition}[Well-founded order]\label{def:b1:induction-and-termination:well-founded}
A \omterm{def:b1:induction-and-termination:strict-order}{strict order} is \emph{well-founded}\index{well-founded order} if there is no infinite strictly descending chain $x_0
> x_1 > x_2 > \cdots$.
\end{definition}

The natural numbers with $<$ are \omterm{def:b1:induction-and-termination:well-founded}{well-founded}; the integers are not ($0 > -1 > -2 > \cdots$), nor the non-negative
rationals ($1 > 1/2 > 1/4 > \cdots$). Proper \omterm{def:b1:induction-and-termination:rooted-tree}{subtrees} of a tree, or the order ``is a proper \omterm{def:b1:propositional-logic:size}{subformula} of'', are
\omterm{def:b1:induction-and-termination:well-founded}{well-founded}, since each step down a tree loses at least one node.

\begin{proposition}[Two definitions agree]\label{prop:b1:induction-and-termination:minimal}
A \omterm{def:b1:induction-and-termination:strict-order}{strict order} is \omterm{def:b1:induction-and-termination:well-founded}{well-founded} if and only if every non-empty subset has a \omterm{def:b1:induction-and-termination:strict-order}{minimal element}.
\end{proposition}

\begin{proof}
If a non-empty set $S$ has no \omterm{def:b1:induction-and-termination:strict-order}{minimal element}, pick $x_0$ in $S$; it is not \omterm{def:b1:induction-and-termination:strict-order}{minimal}, so some $x_1 \in S$ is smaller;
and so on, which builds an infinite descending chain. Conversely, an infinite descending chain is a non-empty set
without \omterm{def:b1:induction-and-termination:strict-order}{minimal element}: each of its elements has a smaller one after it.
\end{proof}

The second form gives \emph{\omterm{def:b1:induction-and-termination:well-founded}{well-founded} induction}: if $P(y)$ follows from $P(x)$ for all $x < y$, then $P$ holds
everywhere, for otherwise the set of counterexamples would have a \omterm{def:b1:induction-and-termination:strict-order}{minimal element} $y$, all of whose smaller elements
satisfy $P$, and so would $y$. Induction on $\mathbb{N}$ and \omterm{def:b1:induction-and-termination:structural}{structural induction} are the cases of $<$ on numbers and
of the ``is a proper \omterm{def:b1:induction-and-termination:rooted-tree}{subtree} of'' order.

\begin{definition}[Product order, lexicographic order]\label{def:b1:induction-and-termination:lex}
On pairs of elements of ordered sets $A$ and $B$, the \emph{product order}\index{product order} sets $(a, b) \leq (a',
b')$ when $a \leq a'$ and $b \leq b'$; the \emph{lexicographic order}\index{lexicographic order} sets $(a, b) < (a', b')$
when $a < a'$, or $a = a'$ and $b < b'$. Both extend to tuples of a fixed length $k$, the lexicographic order comparing
the first component where two tuples differ.
\end{definition}

\begin{proposition}[Lexicographic order is well-founded]\label{prop:b1:induction-and-termination:lex}
If $<_A$ and $<_B$ are \omterm{def:b1:induction-and-termination:well-founded}{well-founded}, the \omterm{def:b1:induction-and-termination:lex}{lexicographic order} on $A \times B$ is \omterm{def:b1:induction-and-termination:well-founded}{well-founded}, and so is the strict
\omterm{def:b1:induction-and-termination:lex}{product order}. By induction, so are both orders on $k$-tuples.
\end{proposition}

\begin{proof}
Take a descending chain $(a_0, b_0) > (a_1, b_1) > \cdots$. Along it, $a_0 \geq a_1 \geq \cdots$, since each step
either decreases $a$ or keeps it. As $A$ is \omterm{def:b1:induction-and-termination:well-founded}{well-founded}, $a$ decreases only finitely many times: from some index on,
$a$ is constant, and then $b$ decreases at every step, which can happen only finitely often in $B$. So the chain is
finite. A strict decrease in the \omterm{def:b1:induction-and-termination:lex}{product order} is a decrease in the \omterm{def:b1:induction-and-termination:lex}{lexicographic order}, so the \omterm{def:b1:induction-and-termination:lex}{product order} has no
infinite chain either. For $k$-tuples, write $A_1 \times \dots \times A_k$ as $A_1 \times (A_2 \times \dots \times A_k)$.
\end{proof}

\begin{omfigure}
\begin{tikzpicture}[x=6mm, y=6mm, >={Stealth[length=2mm]}]
\foreach \i in {0,...,5} {\foreach \j in {0,...,7} {\fill[black!25] (\i,\j) circle (1.2pt);}}
\draw[->] (-0.5,0) -- (6,0) node[right, font=\scriptsize] {$i$};
\draw[->] (0,-0.5) -- (0,8) node[above, font=\scriptsize] {$j$};
\foreach \a/\b/\c/\d in {4/1/4/0, 4/0/3/7, 3/7/3/4, 3/4/3/0, 3/0/2/6, 2/6/2/0, 2/0/1/3, 1/3/1/0, 1/0/0/2, 0/2/0/0} {
  \draw[->, thick, omThm] (\a,\b) -- (\c,\d);
}
\fill[omThm] (4,1) circle (2.5pt);
\node[font=\scriptsize, anchor=west, align=left] at (6.6,4.5) {each step lowers $j$ with $i$ fixed,\\or lowers $i$ and
sets $j$ to any value;\\$i$ can drop only 4 times, and between\\two drops $j$ only finitely often:\\the chain stops at
$(0, 0)$};
\end{tikzpicture}
\par\smallskip{\small A descending chain in the \omterm{def:b1:induction-and-termination:lex}{lexicographic order} on $\mathbb{N}^2$ from $(4, 1)$. A step may jump
upwards in $j$, without bound, but only by lowering $i$: the chain is finite, although no bound on its length can be
given in advance.}
\end{omfigure}

The \omterm{def:b1:induction-and-termination:lex}{lexicographic order} on $\mathbb{N}^2$ shows why well-foundedness is more than ``the values stay positive'': the
pair $(1, 0)$ has infinitely many smaller elements $(0, j)$, and no \omterm{def:b1:induction-and-termination:strict-order}{immediate predecessor}, yet every descending chain is
finite. A directed graph whose arcs all go from a larger to a smaller element has no cycle; the converse view, the
order of reachability in a directed acyclic graph, returns with graphs in \cref{ch:b1:graph-traversals}.

\section{Termination by a variant}\label{sec:b1:induction-and-termination:variant}

\begin{definition}[Termination, variant]\label{def:b1:induction-and-termination:variant}
A loop or a recursive function \emph{terminates}\index{termination} if every execution reaches its end after finitely
many steps. A \emph{variant}\index{variant (termination)} is a function from the states of the loop (or the arguments
of the function) to a set with a \omterm{def:b1:induction-and-termination:well-founded}{well-founded order} that strictly decreases at each iteration (or at each recursive
call). A loop with a variant terminates: otherwise the values of the variant would form an infinite descending chain.
\end{definition}

Euclid's algorithm, which replaces $(a, b)$ by $(b, a \bmod b)$ while $b > 0$, has the \omterm{def:b1:induction-and-termination:variant}{variant} $b$ in $\mathbb{N}$:
$a \bmod b < b$. A \omterm{def:b1:induction-and-termination:variant}{variant} need not be a number. The loop below has none in $\mathbb{N}$, since it sometimes increases
$j$, but the pair $(i, j)$ decreases in the \omterm{def:b1:induction-and-termination:lex}{lexicographic order} at every step.

\omcode{code/bachelor-1/11-induction-and-termination/listings/lex_loop.py}{3}{12}{A loop with a lexicographic
\omterm{def:b1:induction-and-termination:variant}{variant}}\label{lst:b1:induction-and-termination:lex-loop}

\omoutput[1-5,14-17]{code/bachelor-1/11-induction-and-termination/outputs/lex_loop.txt}{The \omterm{def:b1:induction-and-termination:variant}{variant} at each step (the
middle steps are left out).}\label{out:b1:induction-and-termination:lex-loop}

Ackermann's function, named after the mathematician Wilhelm Ackermann, \omterm{def:b1:induction-and-termination:variant}{terminates} for the same reason, although it
calls itself on an argument that is the result of another call.

\omcode{code/bachelor-1/11-induction-and-termination/listings/ackermann.py}{9}{16}{Ackermann's
function}\label{lst:b1:induction-and-termination:ackermann}

\omoutput{code/bachelor-1/11-induction-and-termination/outputs/ackermann.txt}{Values, and calls in parentheses: the third
row already doubles at each step.}\label{out:b1:induction-and-termination:ackermann}

\begin{proposition}[Ackermann's function terminates]\label{prop:b1:induction-and-termination:ackermann}
For all natural numbers $m$ and $n$, the call $\mathit{ack}(m, n)$ \omterm{def:b1:induction-and-termination:variant}{terminates}.
\end{proposition}

\begin{proof}
By \omterm{def:b1:induction-and-termination:well-founded}{well-founded} induction on $(m, n)$ in the \omterm{def:b1:induction-and-termination:lex}{lexicographic order} of $\mathbb{N}^2$. Suppose every call on a smaller
pair \omterm{def:b1:induction-and-termination:variant}{terminates}. If $m = 0$, there is no call. If $n = 0$, the call is on $(m - 1, 1) < (m, 0)$. Otherwise the inner call
is on $(m, n - 1) < (m, n)$ and \omterm{def:b1:induction-and-termination:variant}{terminates} with some value $v$; the outer call is on $(m - 1, v) < (m, n)$, whatever
$v$ is, and \omterm{def:b1:induction-and-termination:variant}{terminates}.
\end{proof}

\omterm{def:b1:induction-and-termination:variant}{Termination} says nothing about speed: $\mathit{ack}(4, 2)$ has 19\,729 decimal digits, and its call count is beyond
reach. The third row is $2^{n + 3} - 3$, which the test checks.

\section{Recursive functions proved correct}\label{sec:b1:induction-and-termination:recursion}

\omcode{code/bachelor-1/11-induction-and-termination/listings/merge.py}{4}{11}{Merging two sorted
lists}\label{lst:b1:induction-and-termination:merge}

\omoutput{code/bachelor-1/11-induction-and-termination/outputs/merge.txt}{Two sorted lists merged, and the base
cases.}\label{out:b1:induction-and-termination:merge}

\begin{proposition}[Merge is correct]\label{prop:b1:induction-and-termination:merge}
If $xs$ and $ys$ are sorted, $\mathit{merge}(xs, ys)$ \omterm{def:b1:induction-and-termination:variant}{terminates} and returns a sorted list holding the elements of
both, with their multiplicities.
\end{proposition}

\begin{proof}
\omterm{def:b1:induction-and-termination:variant}{Termination}: $\text{len}(xs) + \text{len}(ys)$ is a \omterm{def:b1:induction-and-termination:variant}{variant}, since each call removes one element. Correctness, by
induction on the same quantity: if one list is empty, the other is returned, sorted. Otherwise, say $xs[0] \leq
ys[0]$. By induction, the recursive call returns a sorted list of the elements of $xs[1{:}]$ and $ys$; each of them is
at least $xs[0]$, since $xs$ is sorted and $ys[0] \geq xs[0]$. Putting $xs[0]$ in front keeps the list sorted, and adds
the one missing element. The other case is symmetric.
\end{proof}

\begin{definition}[Left-child right-sibling representation]\label{def:b1:induction-and-termination:lcrs}
The \emph{left-child right-sibling representation}\index{left-child right-sibling representation} stores a tree of
any arity as a \omterm{def:b1:induction-and-termination:binary-tree}{binary tree} with the same nodes: the left \omterm{def:b1:induction-and-termination:rooted-tree}{subtree} of a node is its first child, and its right \omterm{def:b1:induction-and-termination:rooted-tree}{subtree}
is its next sibling.
\end{definition}

\omcode{code/bachelor-1/11-induction-and-termination/listings/lcrs.py}{9}{32}{From any arity to binary, and
back}\label{lst:b1:induction-and-termination:lcrs}

\omoutput{code/bachelor-1/11-induction-and-termination/outputs/lcrs.txt}{The tree of the figure below as a \omterm{def:b1:induction-and-termination:binary-tree}{binary tree}
(\code{.} is the empty tree), and the round trip.}\label{out:b1:induction-and-termination:lcrs}

\begin{omfigure}
\begin{tikzpicture}[n/.style={circle, draw, fill=omCode!10, minimum size=5mm, inner sep=0pt, font=\scriptsize},
                    >={Stealth[length=1.6mm]}]
\begin{scope}
\node[n] (1) at (1.8,0) {1};
\foreach \k/\x in {2/0, 3/1.2, 4/2.4, 5/3.6} {\node[n] (\k) at (\x,-1.1) {\k}; \draw (1) -- (\k);}
\node[n] (6) at (0.8,-2.2) {6}; \node[n] (7) at (1.6,-2.2) {7}; \node[n] (8) at (3.6,-2.2) {8};
\draw (3) -- (6); \draw (3) -- (7); \draw (5) -- (8);
\node[font=\scriptsize] at (1.8,-3.0) {a tree of arity 4};
\end{scope}
\begin{scope}[xshift=6cm]
\node[n] (b1) at (0,0) {1};
\node[n] (b2) at (-0.6,-0.9) {2}; \draw[->, thick] (b1) -- (b2);
\node[n] (b3) at (0.3,-1.6) {3}; \draw[->, dashed, omThm] (b2) -- (b3);
\node[n] (b6) at (-0.3,-2.4) {6}; \draw[->, thick] (b3) -- (b6);
\node[n] (b7) at (0.5,-3.1) {7}; \draw[->, dashed, omThm] (b6) -- (b7);
\node[n] (b4) at (1.2,-2.3) {4}; \draw[->, dashed, omThm] (b3) -- (b4);
\node[n] (b5) at (2.1,-3.0) {5}; \draw[->, dashed, omThm] (b4) -- (b5);
\node[n] (b8) at (1.6,-3.8) {8}; \draw[->, thick] (b5) -- (b8);
\node[font=\scriptsize, align=left, anchor=west] at (2.6,-0.9) {solid: left = first child\\dashed: right = next
sibling};
\end{scope}
\end{tikzpicture}
\par\smallskip{\small A tree of arity 4 and its left-child right-sibling \omterm{def:b1:induction-and-termination:binary-tree}{binary tree}. The same eight nodes; the
\omterm{def:b1:induction-and-termination:rooted-tree}{children} of a node become a chain of right links starting at its left child.}
\end{omfigure}

The two conversions are inverse to each other, which is proved by \omterm{def:b1:induction-and-termination:structural}{structural induction} on trees (and on lists of
trees, the two definitions depending on each other); the test checks it on random trees. \omterm{def:b1:c-functions-and-stack:mutual}{Mutually recursive} functions,
such as the even and odd test of \cref{ch:b1:c-functions-and-stack}, terminate by a common \omterm{def:b1:induction-and-termination:variant}{variant}: here $n$ decreases
at every call, whichever function makes it.

\omcode{code/bachelor-1/11-induction-and-termination/listings/even_odd.py}{4}{13}{Even and odd, in
Python}\label{lst:b1:induction-and-termination:even-odd}

\omoutput{code/bachelor-1/11-induction-and-termination/outputs/even_odd.txt}{The even numbers below 12, and seven is
odd.}\label{out:b1:induction-and-termination:even-odd}

\section{Exercises}\label{sec:b1:induction-and-termination:exercises}

\begin{exercise}[$\star$]\label{exo:b1:induction-and-termination:1}
Give the \omterm{def:b1:induction-and-termination:height}{size} and \omterm{def:b1:induction-and-termination:height}{height} of: the empty tree; a single node; a node whose left \omterm{def:b1:induction-and-termination:rooted-tree}{subtree} is a single node and right
\omterm{def:b1:induction-and-termination:rooted-tree}{subtree} empty; and the tree \code{t} of \cref{out:b1:induction-and-termination:inductive}, \code{Node(Node(None, 2,
None), 5, Node(Node(None, 7, None), 9, None))}.
\end{exercise}

\begin{exercise}[$\star$]\label{exo:b1:induction-and-termination:2}
Euclid's algorithm runs on $(252, 105)$. Give the successive values of the \omterm{def:b1:induction-and-termination:variant}{variant} $b$, and the result.
\end{exercise}

\begin{exercise}[$\star$]\label{exo:b1:induction-and-termination:3}
Predict the output of the program.
\omcode{code/bachelor-1/11-induction-and-termination/listings/exo_lcrs.py}{3}{6}{Two small trees converted}
\end{exercise}

\begin{exercise}[$\star$]\label{exo:b1:induction-and-termination:4}
Which \omterm{def:b1:induction-and-termination:strict-order}{strict orders} are \omterm{def:b1:induction-and-termination:well-founded}{well-founded}? (a) $<$ on the integers greater than $-10$; (b) ``divides and is different'' on
positive integers; (c) ``is a proper prefix of'' on finite words; (d) $<$ on the real numbers of $[0, 1]$.
\end{exercise}

\begin{exercise}[$\star\star$]\label{exo:b1:induction-and-termination:5}
A student claims that $n$ is a \omterm{def:b1:induction-and-termination:variant}{variant} of \code{steps}. Show that it is not, prove that the function \omterm{def:b1:induction-and-termination:variant}{terminates} for
all natural numbers $n$ and $\mathit{step} \in \{0, 1\}$ with a lexicographic \omterm{def:b1:induction-and-termination:variant}{variant}, and give its result.
\omcode{code/bachelor-1/11-induction-and-termination/listings/exo_variant.py}{4}{13}{A recursion with two kinds of
call}
\end{exercise}

\begin{exercise}[$\star\star$]\label{exo:b1:induction-and-termination:6}
A library defines the \omterm{def:b1:induction-and-termination:height}{height} of the empty tree as 0. Which statements of
\cref{prop:b1:induction-and-termination:bounds} fail with that convention, and how must they be rewritten?
\omcode{code/bachelor-1/11-induction-and-termination/listings/exo_height0.py}{6}{14}{Another \omterm{def:b1:induction-and-termination:height}{height}}
\end{exercise}

\begin{exercise}[$\star\star$]\label{exo:b1:induction-and-termination:7}
Write \code{count_paths(t)}, the number of paths from the \omterm{def:b1:induction-and-termination:rooted-tree}{root} of a \omterm{def:b1:induction-and-termination:binary-tree}{binary tree} to a \omterm{def:b1:induction-and-termination:rooted-tree}{leaf}, by recursion on the tree.
Prove by \omterm{def:b1:induction-and-termination:structural}{structural induction} that it equals the number of leaves, and give its value on a \omterm{def:b1:induction-and-termination:binary-tree}{full} tree with $i$ \omterm{def:b1:induction-and-termination:rooted-tree}{internal
nodes}.
\end{exercise}

\begin{exercise}[$\star\star$]\label{exo:b1:induction-and-termination:8}
With $\mathit{flatten}(\text{empty}) = [\,]$ and $\mathit{flatten}((\ell, k, r)) = \mathit{flatten}(\ell) + [k] +
\mathit{flatten}(r)$, prove $\text{len}(\mathit{flatten}(t)) = \text{size}(t)$ by \omterm{def:b1:induction-and-termination:structural}{structural induction}.
\end{exercise}

\begin{exercise}[$\star\star$]\label{exo:b1:induction-and-termination:9}
Prove that a non-empty \omterm{def:b1:induction-and-termination:binary-tree}{binary tree} of \omterm{def:b1:induction-and-termination:height}{height} $h$ has at most $2^h$ leaves. Which trees reach the bound?
\end{exercise}

\begin{exercise}[$\star\star\star$]\label{exo:b1:induction-and-termination:10}
On the finite words over a finite alphabet, the \emph{shortlex} order compares lengths first and, between words of
the same length, uses the dictionary order. Prove that it is \omterm{def:b1:induction-and-termination:well-founded}{well-founded}. Show that the dictionary order alone, on the
words over $\{a, b\}$, is not.
\end{exercise}

\begin{exercise}[$\star\star\star$]\label{exo:b1:induction-and-termination:11}
In the sorting game, a move swaps two adjacent numbers of a list that are out of order. Prove that every game ends, with
a \omterm{def:b1:induction-and-termination:variant}{variant}, that it ends on the sorted list, and that the number of moves does not depend on the choices. Write the
game and test the claims.
\end{exercise}

\begin{exercise}[$\star\star\star$]\label{exo:b1:induction-and-termination:12}
McCarthy's 91 function is $f(n) = n - 10$ if $n > 100$, and $f(n) = f(f(n + 11))$ otherwise. Prove that $f(n) = 91$ for
every integer $n \leq 100$, by \omterm{def:b1:induction-and-termination:well-founded}{well-founded} induction on $101 - n$ (with a separate argument for $90 \leq n \leq 100$),
and deduce that $f$ \omterm{def:b1:induction-and-termination:variant}{terminates} on every integer.
\end{exercise}

\section{Problem: The game that always ends}\label{sec:b1:induction-and-termination:problem}

\begin{problem}[{Weekend problem --- The game that always ends}]\label{pb:b1:induction-and-termination:1}
A position is a list of natural numbers. A move picks an entry $k > 0$ and replaces it by two entries $k - 1$; the play
ends when every entry is 0. From $(2)$ one may play $(2) \to (1, 1) \to (0, 0, 1) \to (0, 0, 0, 0)$, or replace the
other 1 first. The program lists the possible lengths of all plays from a few positions, and counts the plays.
\omcode{code/bachelor-1/11-induction-and-termination/listings/game.py}{6}{30}{The game, its plays and their
lengths}\label{lst:b1:induction-and-termination:game}
\omoutput{code/bachelor-1/11-induction-and-termination/outputs/game.txt}{One length per position, whatever the
choices, among up to 21\,964\,800 different plays.}\label{out:b1:induction-and-termination:game}

\needspace{4\baselineskip}
\medskip\noindent\textbf{Part I --- Positions and plays as data.}
\begin{enumerate}
\item Define the set of plays inductively: a play from a position $s$ is either $s$ alone, if no move is possible, or
$s$ followed by a play from a position one move away.
\item List the two plays from $(2)$ and the two from $(1, 1)$.
\item Why does \code{plays((4,))} equal \code{plays((3, 3))}?
\item Explain the recursion of \code{lengths}, and why its base case returns $\{0\}$.
\item Why is \code{@cache} needed to count 21\,964\,800 plays?
\item Does the order of the entries in a position matter for the lengths? For the number of plays?
\end{enumerate}

\needspace{4\baselineskip}
\medskip\noindent\textbf{Part II --- Termination.}
\begin{enumerate}[resume]
\item Show that the sum of the entries is not a \omterm{def:b1:induction-and-termination:variant}{variant}, nor the number of entries, nor the largest entry.
\item Let $K$ be the largest entry of the starting position. Show that no entry ever exceeds $K$.
\item For a position $s$, let $c(s) = (c_K, c_{K-1}, \dots, c_0)$, where $c_j$ is the number of entries equal to $j$.
Show that a move decreases $c(s)$ in the \omterm{def:b1:induction-and-termination:lex}{lexicographic order} of $\mathbb{N}^{K+1}$.
\item Deduce from \cref{prop:b1:induction-and-termination:lex} that every play is finite.
\item Is $c$ a \omterm{def:b1:induction-and-termination:variant}{variant} in the \omterm{def:b1:induction-and-termination:lex}{product order}? Why does that not matter?
\item Could a play from $(3, 1, 0)$ be longer than 1000 moves? Answer from Part II alone, then from the table.
\end{enumerate}

\needspace{4\baselineskip}
\medskip\noindent\textbf{Part III --- A single entry.}
\begin{enumerate}[resume]
\item Let $M(k)$ be the number of moves of a play from $(k)$, assuming it does not depend on the choices. Show
$M(0) = 0$ and $M(k) = 1 + 2M(k - 1)$.
\item Deduce $M(k) = 2^k - 1$ by induction, and compare with the table.
\item In a play from $(k)$, the two copies of $k - 1$ evolve independently. Make this precise: a play from a position
is an interleaving of plays from its entries.
\item Deduce that a play from $(a_1, \dots, a_n)$ makes $\sum_i (2^{a_i} - 1)$ moves, if each single-entry play has a
fixed length.
\item Where did Part III assume what it set out to prove? Why does the interleaving argument still need the
single-entry result as an induction hypothesis?
\item Give the number of plays from $(2)$, $(3)$ and $(1, 2)$ from the table, and check $(2)$ and $(1, 2)$ by hand.
\end{enumerate}

\needspace{4\baselineskip}
\medskip\noindent\textbf{Part IV --- A potential.}
\begin{enumerate}[resume]
\item Let $\Phi(s) = \sum_i (2^{a_i} - 1)$. Compute the change of $\Phi$ in a move that replaces $k$ by two copies of $k -
1$.
\item Show that $\Phi(s) \geq 0$, with equality exactly when the play has ended.
\item Deduce that every play from $s$ has exactly $\Phi(s)$ moves.
\item Compare this proof with that of Part II: which one gives \omterm{def:b1:induction-and-termination:variant}{termination}, which one the length, and why is
$\Phi$ a \omterm{def:b1:induction-and-termination:variant}{variant} in $\mathbb{N}$ although the sum of entries was not?
\item The program checks the result on every position of at most three entries below 4. What does that check add to
the proof, and what does the proof add to the check?
\item State the named result: \emph{every play from $(a_1, \dots, a_n)$ has the same length $\sum_i (2^{a_i} - 1)$}.
\end{enumerate}
\end{problem}
